\chapter*{List of Abbreviations and Symbols}

\addcontentsline{toc}{chapter}{List of Abbreviations and Symbols}

\renewcommand{\arraystretch}{1.3}

\subsection*{Abbreviations and Acronyms}

\begin{longtable}{|c|p{10.5cm}|}
    \hline
    \textbf{Term} & \textbf{Meaning} \\
    \hline
    \endhead
    QCD & Quantum Chromodynamics \\ \hline
    LO / NLO / NNLO & Leading / Next-to-Leading / Next-to-Next-to-Leading Order \\ \hline
    IR / UV & Infrared / Ultraviolet \\ \hline
    MI & Master Integral(s) \\ \hline
    RR / RV / VV & double-Real, Real-Virtual and double-Virtual NNLO contributions \\ \hline
    AP & Altarelli--Parisi (splitting function) \\ \hline
    CDR & Conventional Dimensional Regularisation \\ \hline
    KLN & Kinoshita--Lee--Nauenberg (theorem) \\ \hline
    SM & Standard Model of Particle Physics \\ \hline
    EFT & Effective Field Theory \\ \hline
    SUSY & Supersymmetry \\ \hline
    MC & Monte Carlo \\ \hline
    IBP & Integration by Parts \\ \hline
\end{longtable}

\subsection*{Antenna and Regularisation Notation}

\begin{longtable}{|c|p{10.5cm}|}
    \hline
    \textbf{Symbol} & \textbf{Meaning} \\
    \hline
    \endhead
    $X_n^l$ & Generic antenna symbol: $n$ is the final-state multiplicity, $l$ the loop order \\ \hline
    $A$, $B$, $C$ & Quark-antiquark antenna families (implemented) \\ \hline
    $D$, $E$ & Quark-gluon antenna families (not implemented in the present work) \\ \hline
    $F$, $G$, $H$ & Gluon-gluon antenna families (not implemented in the present work) \\ \hline
    $X$ (undecorated) & Leading-colour component (from the raw ``Lead'' interference term) \\ \hline
    $\tilde{X}$ & Subleading-colour component (raw ``Sublead'' interference term) \\ \hline
    $\hat{X}$ & Quark-loop component (raw ``QL'' interference term) \\ \hline
    $\breve{X}$ & Double-fermion-loop self-interference component (NNLO $A_2^2$ only) \\ \hline
    $S_\epsilon$ & $(4\pi)^{-\epsilon}e^{\epsilon\gamma_E}$, the $\overline{\text{MS}}$ normalisation factor \\ \hline
    $\Lambda_l$ & $(8\pi^2)^l$, loop-order normalisation \\ \hline
    $G_k$ & $(8\pi^2)^k$, perturbative-order scale normalisation \\ \hline
    $C(\epsilon,k)$ & $G_kS_\epsilon$, combined normalisation factor \\ \hline
    $\mathcal N_X$ & Build-stage colour and coupling normalisation \\ \hline
    $\mathcal N(\epsilon,k)$ & Full integration normalisation, $C(\epsilon,k)/\Phi_2$ \\ \hline
    $\mathcal T$-objects & Process-wide colour- and kinematics-conscious objects assembled from antennae, used in the $R$-ratio construction \\ \hline
    $\Phi_n$ & $d$-dimensional $n$-body phase-space volume \\ \hline
\end{longtable}

\subsection*{General Kinematics}

\begin{longtable}{|c|p{10.5cm}|}
    \hline
    \textbf{Symbol} & \textbf{Meaning} \\
    \hline
    \endhead
    $\alpha_s$ & Strong coupling constant \\ \hline
    $\epsilon$, $d=4-2\epsilon$ & Dimensional-regularisation parameter; spacetime dimension \\ \hline
    $s_{ij}$ & Mandelstam invariant, $2p_i\cdot p_j$ \\ \hline
    $\mu$ & Renormalisation/regularisation scale \\ \hline
    $q^2$ & Hard-process invariant mass squared \\ \hline
    $N$ & Number of colours ($=3$) \\ \hline
    $N_f$ & Number of active quark flavours \\ \hline
    $C_F$, $C_A$ & Quadratic Casimir invariants (fundamental / adjoint representation) \\ \hline
    $g_s$ & Strong coupling, $g_s=\sqrt{4\pi\alpha_s}$ \\ \hline
\end{longtable}
